% Source: https://artofproblemsolving.com/community/c3h1221842_the_right_circumcircle_problem
% Stable identity: aops-topic-1221842-q01
The sides are $\dfrac{5}{2}$ times $3,4,5$, so the triangle is right-angled and its hypotenuse $12\dfrac{1}{2}$ is a diameter of the circumcircle. The radius is $\dfrac{25}{4}$: $\boxed{\text{(C)}}$.

\begin{asy}
usepackage("amsmath");
size(3cm);
defaultpen(fontsize(10pt));
pair A=(0,0),B=(10,0),C=(10,7.5); draw(A--B--C--cycle,linewidth(1)); draw(circle((A+C)/2,6.25),blue); dot((A+C)/2);
draw((9.6,0)--(9.6,0.4)--(10,0.4)); label("$10$",(5,0),S); label("$7\dfrac{1}{2}$",(10,3.75),E); label("$12\dfrac{1}{2}$",(5,3.75),NW);
label("\mbox{}",(min(currentpicture,user=true)+max(currentpicture,user=true))/2,black); // keeps the text after the figure black when LaTeX typesets the labels inline
\end{asy}
